\section*{Bab \ref{ch:b3:colloids} --- Antarmuka, Surfaktan, dan Koloid}

\begin{solution}{exo:b3:colloids:1}
$r = \qty{0.50}{\micro m}$: $\Delta p = 2 \times 0.0720/\num{5.0e-7} = \qty{2.9e5}{Pa}$; di dalamnya, $1.0 + 2.9 = \qty{3.9}{bar}$.
\end{solution}

\begin{solution}{exo:b3:colloids:2}
$\ln(p/p^*) = 2 \times 0.0720 \times \num{1.807e-5}/(\num{1.0e-8} \times 8.314 \times 298.15) = 0.105$; $p/p^* = 1.11$.
\end{solution}

\begin{solution}{exo:b3:colloids:3}
$\cos\theta = (40 - 20)/72 = 0.278$, $\theta = \ang{74}$: di bawah $90^\circ$, cairan itu \omterm{def:b3:colloids:wetting}{membasahi} padatan
sebagian.
\end{solution}

\begin{solution}{exo:b3:colloids:4}
NaCl: $I = \qty{10}{mol.m^{-3}}$, $\kappa^{-1} = 0.96\sqrt{100/10} = \qty{3.0}{nm}$. \ce{MgSO4}: $I = \qty{40}{mol.m^{-3}}$, \qty{1.5}{nm}.
\end{solution}

\begin{solution}{exo:b3:colloids:5}
$\Gamma = \num{12.6e-3}/(2 \times 8.314 \times 298.15) = \qty{2.54e-6}{mol.m^{-2}}$; luas $1/(\Gamma N_A) =
\qty{0.65}{nm^2}$ per molekul.
\end{solution}

\begin{solution}{exo:b3:colloids:6}
Di bawah CMC kemiringannya $-12.9$ \unit{mN/m} per satuan $\log c$ (dari $-5.0$ sampai $-4.0$);
dataran berada pada 33.4. Garis itu mencapainya pada $\log c = -4.0 + (39.1 - 33.4)/12.9 = -3.56$: CMC
$\approx \qty{2.8e-4}{mol/L}$.
\end{solution}

\begin{solution}{exo:b3:colloids:7}
SDS: $P = 0.350/(0.60 \times 1.67) = 0.35$, di perbatasan antara bola dan
silinder pendek: \omterm{def:b3:colloids:micelle}{misel} bulat di dekat CMC. Lipid: $P = 0.700/(0.65 \times 1.67) = 0.64$:
lapisan ganda, jadi vesikel.
\end{solution}

\begin{solution}{exo:b3:colloids:8}
\ce{CaCl2}: $60/64 = \qty{0.94}{mM}$ \ce{Ca^{2+}}; \ce{AlCl3}: $60/729 = \qty{0.082}{mM}$ \ce{Al^{3+}}. \ce{Na3PO4}
mengoagulasi melalui \ce{Na+}-nya (pada sekitar \qty{20}{mM} garam); fosfatnya, sebuah
ion sejenis, bahkan dapat teradsorpsi dan menaikkan muatan negatifnya.
\end{solution}

\begin{solution}{exo:b3:colloids:9}
Bagian hidrofiliknya $6 \times 44.0 + 17.0 = \qty{281.0}{g/mol}$ dari \qty{450.0}{g/mol} (massa atom buku ini):
HLB $= 20 \times 281.0/450.0 = 12.5$, sebuah pengemulsi minyak dalam air.
\end{solution}

\begin{solution}{exo:b3:colloids:10}
Minyak tetesan kecil lebih larut dalam fase kontinu (dengan faktor
$\exp(2\gamma V_{\mathrm m}/rRT)$, yang eksponennya sepuluh kali lebih besar untuk \qty{50}{nm} daripada untuk \qty{500}{nm}):
minyak itu berdifusi ke tetesan besar, yang tumbuh sementara tetesan kecil
lenyap. \omterm{def:b3:colloids:ripening}{Pematangan Ostwald} ini membuat \omterm{def:b3:colloids:colloid}{emulsi} makin kasar bahkan jika tidak
ada dua tetesan pun yang pernah bergabung.
\end{solution}

\begin{solution}{exo:b3:colloids:11}
Meniskusnya berupa tudung bola berjari-jari $r/\cos\theta$: cairan tepat di bawahnya bertekanan
lebih rendah sebesar $2\gamma\cos\theta/r$, dan kolomnya naik sampai $\rho gh$ mengimbanginya. Untuk air,
$h = 2 \times 0.072/(997 \times 9.81 \times \num{1.0e-4}) = \qty{0.15}{m}$.
\end{solution}

\begin{solution}{exo:b3:colloids:12}
Pada \qty{100}{mM} \omterm{def:b3:electrode-kinetics:double-layer}{panjang Debye} kurang dari \qty{1}{nm} dan tolakan lenyap sebelum
tarikan: tidak ada penghalang yang tersisa (melampaui kurva \qty{50}{mM}). Polimer bebas tidak
bekerja di permukaan; lapisan polimer yang tercangkok menolak saat bersentuhan
berapa pun kadar garamnya, karena entropi yang hilang ketika rantai-rantainya
bertumpang tindih.
\end{solution}

\begin{solution}{pb:b3:colloids:1}
\textbf{1.} Substitusi dalam kisi kristal (\ce{Al^{3+}} menggantikan \ce{Si^{4+}}, \ce{Mg^{2+}} menggantikan \ce{Al^{3+}}) menyisakan
muatan negatif yang permanen, yang diimbangi oleh kation-kation yang dapat dipertukarkan.
\textbf{2.} $I = \frac12(1.0 + 1.0) = \qty{1.0}{mM}$.
\textbf{3.} $\kappa^{-1} = \qty{9.6}{nm}$.
\textbf{4.} $a\kappa = 10$: lapisan rangkap tipis dibandingkan dengan partikelnya, yang
merupakan syarat pendekatan Derjaguin.
\textbf{5.} $v = 2 \times 1650 \times 9.81 \times (\num{1.0e-7})^2/(9 \times \num{0.89e-3}) = \qty{4.0e-8}{m/s}$, sekitar \qty{3.5}{mm} per hari.
\textbf{6.} Lapisan-lapisan difusnya saling menolak sebelum tarikan van der Waals mengambil alih.
\textbf{7.} $V = 2\pi\varepsilon_r\varepsilon_0a\psi^2\ln(1 + \eu^{-\kappa h}) - Aa/12h$.
\textbf{8.} Sekitar $39\,kT$.
\textbf{9.} Fraksi berorde $\eu^{-39} \approx 10^{-17}$ dari tumbukan-tumbukannya: praktis tidak pernah.
\textbf{10.} Penghalang itu lenyap: setiap tumbukan berujung pada sentuhan.
\textbf{11.} \omterm{def:b3:electrode-kinetics:double-layer}{Panjang Debye} yang lebih pendek membatasi tolakan pada celah yang kecil,
tempat tarikan, yang tumbuh sebagai $1/h$, mendominasi.
\textbf{12.} CCC sebanding dengan $z^{-6}$, yaitu pangkat enam negatif muatan ion pengimbang.
\textbf{13.} $50/64 = \qty{0.78}{mM}$.
\textbf{14.} $50/729 = \qty{0.069}{mM}$.
\textbf{15.} Kation-kation itu adalah ion pengimbang partikel yang bermuatan
negatif: kation-kation itu menumpuk di dalam lapisan rangkap dan
menyaringnya.
\textbf{16.} Sulfat adalah ion sejenis, yang ditolak oleh partikel-partikel itu.
\textbf{17.} \qty{0.069}{mol} \ce{Al^{3+}} per meter kubik.
\textbf{18.} \qty{0.034}{mol} $\ce{Al2(SO4)3}$.
\textbf{19.} $0.0343 \times 342.3 = \qty{11.7}{g}$.
\textbf{20.} \qty{11.7}{kg} per jam.
\textbf{21.} $\ce{Al^3+ + 3 HCO3- -> Al(OH)3 + 3 CO2}$.
\textbf{22.} $3 \times 0.069 = \qty{0.21}{mM}$ dari \qty{1.0}{mM}: seperlimanya. Hidrogenkarbonat
yang tersisa dan \ce{CO2} terlarut menyangga air itu: pH-nya hanya turun
sedikit.
\textbf{23.} Spesies aluminium yang muatannya tinggi teradsorpsi dan membalik muatan
partikel, yang lalu saling menolak lagi.
\textbf{24.} \textbf{Sekitar \qty{12}{g} aluminium sulfat per meter kubik} (\qty{11.7}{g}), sebelum tambahan
yang diperlukan untuk hidrolisis dalam praktik.
\end{solution}
